\section{Introduction}
\label{sec:introduction}

Sentiment analysis assigns a polarity to a piece of text. On social media this is harder than
on reviews or news: comments are short, informal and full of slang, emojis, sarcasm and
negations, and a single word such as ``not'' can reverse the meaning of the sentence. GoEmotions
\cite{demszky2020goemotions} is a corpus of English Reddit comments annotated with 27 emotions
plus neutral, together with an official grouping of the emotions into positive, negative and
ambiguous. In this work that grouping is used to turn the emotion labels into a binary
sentiment task (positive vs.\ negative).

The main objective is to compare three neural architectures that differ in how they handle
the order of the words and how much they remember: a dense network over word embeddings, which
has no memory of the sequence; a vanilla recurrent network (RNN) \cite{elman1990finding}, which
keeps a hidden state but suffers from vanishing gradients \cite{bengio1994learning}; and a Long
Short-Term Memory network (LSTM) \cite{hochreiter1997long}, whose gates control what is
written, kept and read from its memory cell. All of them share the same preprocessing and
vectorization, and they are compared with a DummyClassifier baseline using accuracy, precision,
recall, F1 and Cohen's kappa \cite{cohen1960coefficient}. As an extension we fine-tune
DistilBERT \cite{sanh2019distilbert}, a smaller version of BERT \cite{devlin2019bert} based on
the attention mechanism of the Transformer \cite{vaswani2017attention}.

A second objective is to relate these models to the theory of computation studied in the
course \cite{sipser2013introduction}. A recurrent network reads its input one symbol at a time
and updates a state, much like the transition function of an automaton; with unbounded
precision RNNs are Turing complete \cite{siegelmann1995computational}, and so are Transformers
\cite{perez2021attention}. Neural Turing Machines \cite{graves2014neural} make the analogy
explicit with an external memory and differentiable read and write heads. In practice, however,
every trained classifier performs a fixed number of operations per input and always halts, so
it decides a language of ``positive comments'' only approximately.

The rest of the report is organized as follows. Section~\ref{sec:data} describes the dataset,
the sentiment mapping and the preprocessing. Section~\ref{sec:models} presents the neural
architectures and their training, and Section~\ref{sec:evaluation} the evaluation results.
Section~\ref{sec:transformer} covers the Transformer extension, Section~\ref{sec:comparison}
compares all models and discusses the relation with Turing machines, and
Section~\ref{sec:conclusions} concludes.
